\subsection{日式环}

定义 1．——设 A 是整诺特环。若 A 在其分式域的每个有限扩张中的整闭包都是有限 A-代数，则称 A 为日式环。

注记．——1) 这等价于说 A 满足如下条件：每个作为 A-代数、在 A 上整且包含于 A 的分式域 K 的某个有限生成扩张中的整环 B，都是有限 A-代数。事实上，B 的分式域 L 是 K 的代数扩张，因而是 K 的有限次扩张（A，V，p. 112，命题 17 的推论 1）。A-代数 B 包含于 A 在 L 中的整闭包内；若后者有限，则 B 也有限。

\begin{enumerate}
\def\labelenumi{\arabic{enumi})}
\setcounter{enumi}{1}
\tightlist
\item
  设 A 是诺特整日式环，S 是 A 的一个不含 0 的乘法子集。分式环 \(\mathbf{S}^{-1}\mathbf{A}\) 是日式环。事实上，令 L 为 A 的分式域的一个有限扩张，令 B 为 A 在 L 中的整闭包；则 \(\mathbf{S}^{-1}\mathbf{A}\) 在 L 中的整闭包为 \(\mathbf{S}^{-1}\mathbf{B}\)（V，§ 1，第 5 号，命题 16），因而是有限的 \(\mathbf{S}^{-1}\mathbf{A}\)-代数。
\end{enumerate}

例．——域上的每个有限型整代数都是日式环（V，§ 3，第 2 号，定理 2）。

命题 1．——设 A 是诺特整环，K 是其分式域。假设对于 K 的每个有限根式扩张 L，A 在 L 中的整闭包都是有限 A-代数。则环 A 是日式的。

令 E 为 K 的有限扩张。令 N 为包含 E 的 K 的有限拟伽罗瓦扩张（A，V，p. 54，推论 1），并令 L 为 N 的 K-自同构群的不变域。于是（A，V，p. 73，命题 13），L 是 K 的根式扩张，而 N 是 L 的可分扩张。因此根据假设，A 在 L 中的整闭包 B 是有限 A-代数；B 在 N 中的整闭包 C 是有限 B-代数（V，§ 1，第 6 号，命题 18 的推论 1），因而是有限 A-代数。A 在 E 中的整闭包包含于 C 内，所以由于 A 是诺特环，它是有限 A-代数。

推论．——假设域 K 是完美的（例如特征为 0）。则 A 是日式环，当且仅当其整闭包是有限 A-代数。

命题 2．——设 B 是诺特整环，A 是 B 的诺特子环，且 B 是有限 A-代数。A 为日式环的充要条件是 B 为日式环。

令 K（分别为 L）表示 A（分别为 B）的分式域。先假设 A 是日式环，并令 M 为 L 的有限扩张。令 C 为 B 在 M 中的整闭包。根据 V，§ 1，第 1 号，命题 6，C 是 A 在 M 中的整闭包；由于 M 是 K 的有限扩张且 A 是日式环，C 是有限 A-代数。更不用说，C 是有限 B-代数。这就证明了 B 是日式环。

反之，假设 B 是日式环，并令 N 为 K 的有限扩张。令 D 为 A 在 N 中的整闭包。令 E 为 L 与 N 在 K 上的复合；由于 B 是日式环，B 在 E 中的整闭包 D' 是有限 B-代数，因而是有限 A-代数；作为 D' 的子模，A-模 D 因而有限生成，这就说明 A 是日式环。

